\chapter{Diarrhea}

\section*{Definition}
Frequent passage of loose or watery stools exceeding usual frequency and consistency.

\section*{When to See a Doctor}
See your doctor if:
\begin{itemize}
  \item Seek prompt care for dehydration signs such as very little urine, dizziness, dry mouth, unusual sleepiness, or few tears in a child.
  \item Seek care for blood or pus in the stool, severe or persistent abdominal pain, high fever, repeated vomiting, recent antibiotics, or diarrhea that is not improving after several days.
  \item Infants, young children, older adults, and people with major medical conditions may need earlier assessment.
\end{itemize}

\section*{How It Develops}
Caused by increased secretion, decreased absorption, altered transit time, or osmotic load; acute vs chronic guides evaluation.

\section*{Causes}
\begin{itemize}
  \item Viral gastroenteritis
  \item Bacterial infection
  \item Antibiotic-associated diarrhea
  \item IBS-D
  \item Inflammatory bowel disease
\end{itemize}

\section*{Related Symptoms}
\begin{itemize}
  \item Abdominal cramps
  \item Urgency
  \item Dehydration
  \item Electrolyte imbalance
  \item Weight loss
\end{itemize}

\section*{Medical Evaluation}
\begin{itemize}
  \item History assesses duration, stool characteristics, exposure, travel, medicines, recent antibiotics, immune status, fever, and blood.
  \item Examination focuses on vital signs, hydration, abdominal findings, and complications.
  \item Stool testing, blood tests, imaging, or endoscopy is reserved for severe, persistent, bloody, inflammatory, immunocompromised, or otherwise concerning presentations.
\end{itemize}

\section*{Treatment Options}
\begin{itemize}
  \item Oral rehydration is the mainstay for uncomplicated acute diarrhea; selected infections require targeted antibiotics and some inflammatory diseases require specialist therapy.
  \item Chronic diarrhea requires evaluation for malabsorption, medication effects, inflammatory disease, endocrine disorders, and functional bowel disorders.
\end{itemize}

\section*{Self-Care and Home Management}
\begin{itemize}
  \item Replace fluid and salts with oral rehydration solution; continue light meals as tolerated rather than fasting.
  \item Avoid alcohol and foods that clearly worsen symptoms.
  \item Loperamide may help some adults with uncomplicated watery diarrhea, but avoid it with fever, bloody stools, suspected invasive infection, or possible C. difficile infection unless a clinician advises otherwise.
  \item Seek care for dehydration, severe pain, blood or black stool, high fever, persistent vomiting, recent antibiotics, or diarrhea lasting more than a few days.
\end{itemize}

\section*{References}
\begin{thebibliography}{99}
\bibitem{diarrhea:ref1} Riddle MS, DuPont HL, Connor BA. ``ACG clinical guideline: diagnosis, treatment, and prevention of acute diarrheal infections in adults.'' Am J Gastroenterol. 2016;111(5):602-622.
\bibitem{diarrhea:ref2} Farthing M, Salam MA, Lindberg G, et al. ``Acute diarrhea in adults and children: a global perspective.'' J Clin Gastroenterol. 2013;47(1):12-20.
\bibitem{diarrhea:ref3} Thielman NM, Guerrant RL. ``Acute infectious diarrhea.'' N Engl J Med. 2004;350(1):38-47.
\bibitem{diarrhea:ref4} Camilleri M, Murray JA. ``Chronic diarrhea.'' N Engl J Med. 2023;389(15):1402-1414.
\bibitem{diarrhea:ref5} Guerrant RL, Van Gilder T, Steiner TS, et al. ``Practice guidelines for the management of infectious diarrhea.'' Clin Infect Dis. 2001;32(3):331-351.
\bibitem{diarrhea:ref6} Steiner TS. ``Approach to the adult with acute diarrhea.'' UpToDate. Wolters Kluwer; 2025.
\end{thebibliography}
